\subsection{Qualitative Interviews}
\subsubsection{Interviewform und Auswahl der Interviewpartner}
Zur Erhebung praxisnaher Informationen werden in dieser Arbeit qualitative leitfadengestützte Experteninterviews durchgeführt. Ziel der Interviews war es, zunächst einen geeigneten Use Case für einen \ac{KI}-Agenten in \ac{D365FnSCM} zu identifizieren und darauf aufbauend konkrete Anforderungen an den \ac{KI}-Agenten in diesem herausgearbeiteten Use Case zu ermitteln.

Insgesamt werden drei Interviews geführt. Das erste Interview dient der Identifikation eines geeigneten Use Cases. Zwei weitere Interviews werden anschließend eingesetzt, um Erwartungen und Anforderungen an den \ac{KI}-Agenten zu erheben. Die Interviewpartner werden aufgrund ihrer fachlichen Expertise und praktischen Erfahrung im Umgang mit \ac{D365FnSCM} und den Nutzern dieser Systeme ausgewählt. Alle Interviewpartner verfügen über langjährige Erfahrung in der Implementierung und Beratung einer \ac{D365FnSCM}-Lösung.

\subsubsection{Durchführung}
Die Interviews werden alle anhand eines Leitfadens durchgeführt. Der Leitfaden für das erste Interview zur Use-Case-Identifikation umfasst offene Fragen, welche darauf abzielen, einen genauen Anwendungsfall für einen \ac{KI}-Agenten in \ac{D365FnSCM} zu identifizieren (siehe \ref{anhang:InterviewleitfadenUseCase}). Der Leitfaden für die beiden Interviews zur Anforderungserhebung enthalten dabei offene Fragen zu den Erwartungen und Anforderungen an den \ac{KI}-Agenten im identifizierten Use Case (siehe \ref{anhang:InterviewleitfadenAnforderungserhebung}).

Die Interviews werden sowohl schriftlich als auch mündlich durchgeführt. Das erste Interview zur Use-Case-Identifikation sowie eines der beiden Interviews zur Anforderungserhebung werden schriftlich geführt, indem den Interviewpartnern der Interviewleitfaden zugesendet wird. Das dritte Interview zur Anforderungserhebung wird mündlich per Videokonferenz vorgenommen. 

\subsubsection{Aufbereitung und Auswertung}
Die schriftlich durchgeführten Interviews liegen bereits in Textform vor und können somit direkt für die Auswertung genutzt werden. Das mündlich durchgeführte Interview wird aufgezeichnet und anschließend transkribiert. Dabei erfolgt eine sinngemäße Transkription mit Fokus auf die inhaltlichen Aussagen der Interviewpartner. Aspekte wie Pausen, Stottern oder andere sprachliche Merkmale werden nicht berücksichtigt, da der Fokus auf der inhaltlichen Analyse liegt.

Die Auswertung der Interviews erfolgt angelehnt an eine qualitative Inhaltsanalyse nach Mayring (2019). Ziel dieses Vorgehens ist es, die Aussagen der Interviewpartner systematisch zu strukturieren und zentrale Inhalte herauszuarbeiten. Hierzu werden relevante Textpassagen anhand eines Kategoriensystems geordnet, das sich an der Forschungsfrage sowie den Zielsetzungen der Arbeit orientiert. Auf diese Weise können wesentliche Aussagen reduziert, vergleichbar gemacht und in Bezug auf die Anforderungen an den \ac{KI}-Agenten interpretiert werden. Das Vorgehen ermöglicht eine nachvollziehbare und strukturierte Auswertung, ohne den inhaltlichen Kontext der Interviews zu verlieren. \cite[\pagef 633ff]{mayring_qualitative_2019} Das entwickelte Kategoriensystem ist in Tabelle~\ref{tab:kategorien} detailliert dokumentiert.
